% =====================================================================
\chapter{Principios de conteo y ordenaciones}
\label{tema:1}
% =====================================================================

Contar bien es el primer paso para calcular probabilidades en espacios finitos. Este tema reúne las reglas básicas del recuento (suma, producto, complemento, biyección y división), el principio de inclusión--exclusión para dos y tres conjuntos y el recuento de ordenaciones: variaciones y permutaciones, con y sin repetición. La presentación sigue el capítulo~14, «Cardinality Rules», de Lehman, Leighton y Meyer~\cite{mcs}, texto del curso 6.042J del MIT; el capítulo~1 de Ross~\cite{ross} (§1.2 y §1.3); la sección~1.4, «How to count», de Blitzstein y Hwang~\cite{bh}, y la sección~3.1 de Grinstead y Snell~\cite{gs}.

\section{Reglas de la suma, del producto y del complemento}

Se recuerda que, para un conjunto finito $A$, $\abs{A}$ denota su número de elementos y que $[n]=\{1,2,\dots,n\}$. Contar un conjunto consiste en hallar $\abs{A}$; casi siempre se hace descomponiendo $A$ en piezas más sencillas o poniéndolo en correspondencia con otro conjunto ya contado.

\begin{proposicion}[Regla de la suma]\label{prop:t1-suma}
Si $A_1,\dots,A_m$ son conjuntos finitos disjuntos dos a dos, entonces
\[
\abs{A_1\cup\dots\cup A_m}=\abs{A_1}+\dots+\abs{A_m}.
\]
\end{proposicion}

\begin{proposicion}[Regla del producto]\label{prop:t1-producto}
Si $A_1,\dots,A_k$ son conjuntos finitos, $\abs{A_1\times\dots\times A_k}=\abs{A_1}\cdots\abs{A_k}$.
\end{proposicion}

En la práctica se usa una versión más flexible, en la que el conjunto de opciones de cada etapa puede depender de las elecciones previas, pero no su número.

\begin{teorema}[Regla del producto generalizada]\label{teo:t1-producto-gen}
Sea $S$ un conjunto de sucesiones $(s_1,\dots,s_k)$ tal que hay $n_1$ valores posibles de $s_1$ y, para cada elección de $s_1,\dots,s_{i-1}$, hay exactamente $n_i$ valores posibles de $s_i$ ($i=2,\dots,k$). Entonces $\abs{S}=n_1n_2\cdots n_k$.
\end{teorema}

\begin{proposicion}[Regla del complemento]\label{prop:t1-complemento}
Si $A\subseteq U$ y $U$ es finito, $\abs{U\setminus A}=\abs{U}-\abs{A}$.
\end{proposicion}

\addtocounter{definicion}{1}%

\section{Reglas de la biyección y de la división}

\begin{proposicion}[Regla de la biyección]\label{prop:t1-biyeccion}
Si existe una biyección $f\colon A\to B$ entre conjuntos finitos, entonces $\abs{A}=\abs{B}$.
\end{proposicion}

La demostración es inmediata: si $A=\{a_1,\dots,a_n\}$ sin repeticiones, entonces $f(a_1),\dots,f(a_n)$ son distintos (inyectividad) y agotan $B$ (sobreyectividad). Esta regla permite contar un conjunto contando otro~\cite[§14.1]{mcs}.

\addtocounter{definicion}{1}%

\begin{proposicion}[Regla de la división]\label{prop:t1-division}
Sea $f\colon A\to B$ una función sobreyectiva entre conjuntos finitos tal que cada $b\in B$ tiene exactamente $d$ preimágenes ($f$ es \emph{$d$ a~1}). Entonces $\abs{A}=d\,\abs{B}$.
\end{proposicion}

\section{Inclusión--exclusión para dos y tres conjuntos}

Cuando los conjuntos no son disjuntos, la suma de sus cardinales cuenta dos veces los elementos comunes y hay que corregirla.

\begin{teorema}[Inclusión--exclusión]\label{teo:t1-ie}
Para conjuntos finitos $A$, $B$ y $C$:
\begin{align*}
\abs{A\cup B}&=\abs{A}+\abs{B}-\abs{A\cap B},\\
\abs{A\cup B\cup C}&=\abs{A}+\abs{B}+\abs{C}-\abs{A\cap B}-\abs{A\cap C}-\abs{B\cap C}+\abs{A\cap B\cap C}.
\end{align*}
\end{teorema}

La versión para $m$ conjuntos se demuestra en el \cref{tema:2}, una vez disponibles los números combinatorios.

\addtocounter{definicion}{1}%

\section{Variaciones y permutaciones}

Se trata ahora de contar \emph{selecciones ordenadas} de elementos de un conjunto de $n$ elementos. La terminología española distingue según se permita o no repetir.

\begin{definicion}[Variaciones y permutaciones]
Sea $A$ un conjunto con $n$ elementos y $k\ge 0$ un entero.
\begin{enumerate}
  \item Una \emph{variación con repetición} de orden $k$ es una sucesión $(a_1,\dots,a_k)$ de elementos de $A$, posiblemente repetidos; equivale a una función $[k]\to A$.
  \item Una \emph{variación sin repetición} de orden $k\le n$ es una sucesión de $k$ elementos \emph{distintos} de $A$; equivale a una función inyectiva $[k]\to A$.
  \item Una \emph{permutación} de $A$ es una variación sin repetición de orden $n$, es decir, una ordenación de todos sus elementos; equivale a una biyección $[n]\to A$.
\end{enumerate}
Se escribe $n!=1\cdot2\cdots n$, con $0!=1$, y $n^{\underline{k}}=n(n-1)\cdots(n-k+1)$ (\emph{factorial descendente}, con $n^{\underline{0}}=1$).
\end{definicion}

\begin{teorema}\label{teo:t1-variaciones}
Si $\abs{A}=n$, el número de variaciones con repetición de orden $k$ es $n^{k}$; el de variaciones sin repetición de orden $k\le n$ es $n^{\underline{k}}=n!/(n-k)!$, y el de permutaciones es $n!$.
\end{teorema}

\addtocounter{definicion}{1}%

\addtocounter{definicion}{1}%

\addtocounter{definicion}{1}%

\addtocounter{definicion}{1}%

